% ============================================================================
% frontmatter/abstract.tex
% ============================================================================
\chapter*{Abstract}
\addcontentsline{toc}{chapter}{Abstract}

Human Activity Recognition (HAR) using wrist-worn inertial measurement units
has demonstrated strong classification performance on curated benchmark datasets.
However, a significant gap persists between research-grade evaluation and
reliable real-world operation: deployed HAR models degrade over time due to
domain shift caused by inter-user variability, sensor placement differences,
and gradual calibration drift, yet no ground-truth activity labels are available
in production to detect or quantify this degradation.

This thesis designs, implements, and evaluates an end-to-end MLOps pipeline
that closes this gap for an anxiety-related behavioural activity recognition
context.
The pipeline comprises 14 composable stages spanning data ingestion,
preprocessing, windowed inference, three-layer monitoring (confidence, temporal
consistency, and statistical drift), automated trigger policy, model retraining,
and safe deployment with canary evaluation and rollback.

The monitoring framework operates exclusively on unlabelled production
predictions, combining mean confidence tracking, activity transition rate
analysis, Population Stability Index, and Kolmogorov--Smirnov tests into a
composite health signal.
A voting-based trigger engine translates monitoring signals into retraining
decisions using multi-metric confirmation to suppress false alarms.
Two adaptation strategies are implemented and compared: Adaptive Batch
Normalisation followed by TENT entropy minimisation (AdaBN+TENT), and
curriculum pseudo-labelling with Elastic Weight Consolidation regularisation.

Empirical evaluation on 26 Garmin wrist sensor recording sessions demonstrates
that the pipeline detects a physical sensor unit mismatch (milliG vs.\ m/$\text{s}^2$)
that reduced cross-user accuracy to 14.5\,\%, identifies the root cause through
automated drift diagnostics, and recovers model reliability through the
adaptation cycle without any manual re-labelling.

The pipeline is fully containerised with Docker, tracked with MLflow and DVC,
and deployed with a FastAPI inference server behind a GitHub Actions CI/CD
workflow.
All source code, configuration, and experimental artefacts are reproducible
from the repository.

\vspace{0.5cm}
\noindent\textbf{Keywords:} human activity recognition, MLOps, domain adaptation,
wearable sensors, drift detection, pseudo-labelling, AdaBN, TENT,
continuous monitoring, CI/CD.

% --------------------------------------------------------------------------
% German abstract (Zusammenfassung) — required by many German universities
% --------------------------------------------------------------------------
\chapter*{Zusammenfassung}
\addcontentsline{toc}{chapter}{Zusammenfassung}

Die Aktivitätserkennung (Human Activity Recognition, HAR) mit handgelenksgebundenen
inertialen Messeinheiten erzielt auf kuratierten Benchmark-Datensätzen hohe
Klassifikationsleistungen.
In der realen Anwendung treten jedoch erhebliche Herausforderungen auf:
Deployed-Modelle degradieren aufgrund von Domänenwechsel, inter-persönlicher
Variabilität und Sensordrift, ohne dass im Produktionsbetrieb Grundwahrheitslabels
zur Erkennung dieser Degradation verfügbar sind.

Diese Masterarbeit entwirft, implementiert und evaluiert eine End-to-End
MLOps-Pipeline, welche diese Lücke für einen auf angstbezogene Verhaltensklassen
ausgerichteten HAR-Kontext schließt.
Die Pipeline umfasst 14 komposierbare Stufen, darunter Datenaufnahme,
Vorverarbeitung, fensterbasierten Inferenz, dreischichtiges Monitoring
(Konfidenz, temporale Konsistenz, statistische Driftmessung), automatisierte
Trigger-Entscheidung, Modell-Retraining sowie sicheres Deployment mit
Canary-Evaluierung und Rollback-Mechanismus.

\cleardoublepage
